\documentclass[12pt,a4paper]{article}

%% ─── Packages essentiels ─────────────────────────────────────────────────────
\usepackage[utf8]{inputenc}
\usepackage[T1]{fontenc}
\usepackage[french]{babel}
\usepackage{geometry}
\usepackage{graphicx}
\usepackage{amsmath}
\usepackage{booktabs}
\usepackage{array}
\usepackage{xcolor}
\usepackage{hyperref}
\usepackage{fancyhdr}
\usepackage{titlesec}
\usepackage{enumitem}
\usepackage{tcolorbox}
\usepackage{tikz}
\usepackage{pgfplots}
\usepackage{pgf-pie}
\usepackage{multirow}
\usepackage{caption}
\usepackage{subcaption}
\usepackage{float}
\usepackage{parskip}
\usepackage{charter}      % police élégante mais lisible

\pgfplotsset{compat=1.18}

%% ─── Mise en page ────────────────────────────────────────────────────────────
\geometry{
  top=2.2cm, bottom=2.2cm,
  left=2.5cm, right=2.5cm,
  headheight=15pt
}

%% ─── Couleurs personnalisées (harmonie cohérente) ───────────────────────────
\definecolor{primary}{RGB}{30, 90, 160}
\definecolor{secondary}{RGB}{70, 130, 180}
\definecolor{accent}{RGB}{220, 80, 50}
\definecolor{success}{RGB}{34, 139, 34}
\definecolor{warning}{RGB}{255, 140, 0}
\definecolor{lightblue}{RGB}{235, 245, 255}
\definecolor{lightgray}{RGB}{248, 248, 248}
\definecolor{darkgray}{RGB}{60, 60, 60}

%% ─── En-tête / Pied de page ──────────────────────────────────────────────────
\pagestyle{fancy}
\fancyhf{}
\fancyhead[L]{\small\color{primary}\textbf{Projet Scraping \& Dataset}}
\fancyhead[R]{\small\color{darkgray}Rapport final — Tous livrables inclus}
\fancyfoot[C]{\small\color{darkgray}--- \thepage\ ---}
\renewcommand{\headrulewidth}{0.5pt}
\renewcommand{\footrulewidth}{0.3pt}

%% ─── Titres de sections ──────────────────────────────────────────────────────
\titleformat{\section}
  {\Large\bfseries\color{primary}}
  {\thesection.}{0.8em}{}
  [\vspace{-4pt}\rule{\linewidth}{1.2pt}\vspace{6pt}]

\titleformat{\subsection}
  {\large\bfseries\color{secondary}}
  {\thesubsection.}{0.6em}{}

%% ─── Boîtes d'encadré ────────────────────────────────────────────────────────
\tcbuselibrary{skins,breakable}
\newtcolorbox{infobox}[1][]{
  colback=lightblue, colframe=primary,
  fonttitle=\bfseries\small, title=#1,
  arc=4pt, boxrule=0.8pt
}
\newtcolorbox{warnbox}[1][]{
  colback=yellow!10, colframe=warning,
  fonttitle=\bfseries\small, title=#1,
  arc=4pt, boxrule=0.8pt
}
\newtcolorbox{ethicbox}[1][]{
  colback=green!5, colframe=success,
  fonttitle=\bfseries\small, title=#1,
  arc=4pt, boxrule=0.8pt
}

%% ═══════════════════════════════════════════════════════════════════════════
\begin{document}
%% ═══════════════════════════════════════════════════════════════════════════

%% ───────────────────────────────────────────────────────────────────────────
%% PAGE DE TITRE (personnalisée)
%% ───────────────────────────────────────────────────────────────────────────
\begin{titlepage}
  \begin{center}
    \vspace*{1.5cm}
    {\color{primary}\rule{\linewidth}{2pt}}
    \vspace{0.6cm}
    {\Huge\bfseries\color{primary}
      Collecte et Préparation\\[0.3em]
      d'un Dataset d'Images\\[0.3em]
      pour Classification d'Oiseaux}
    \vspace{0.4cm}
    {\color{primary}\rule{\linewidth}{2pt}}
    \vspace{1.2cm}

    \begin{tcolorbox}[colback=lightblue, colframe=primary, arc=6pt,
                      boxrule=1pt, width=0.85\linewidth]
      \centering
      {\large\bfseries Rapport complet}\\[0.4em]
      {\normalsize Méthodologie, statistiques, limites, éthique}
    \end{tcolorbox}

    \vspace{2cm}
    \begin{tabular}{rl}
      \textbf{Repository :} & \href{https://github.com/souhayle98/projet_scraping2}
                              {\color{secondary}\texttt{github.com/souhayle98/projet\_scraping2}} \\[0.4em]
      \textbf{Branches :}   & \texttt{main} (production) \quad \texttt{dev} (développement) \\[0.4em]
      \textbf{Date :}       & \today \\[0.4em]
    \end{tabular}

    \vfill
    \begin{tikzpicture}[node distance=0pt]
      \node[draw=primary, fill=lightblue, rounded corners=4pt,
            minimum width=2.1cm, minimum height=0.75cm,
            font=\small\bfseries, text=primary] (A) at (0,0) {Scraping};
      \node[draw=primary, fill=lightblue, rounded corners=4pt,
            minimum width=2.1cm, minimum height=0.75cm,
            font=\small\bfseries, text=primary] (B) at (2.6,0) {Nettoyage};
      \node[draw=primary, fill=lightblue, rounded corners=4pt,
            minimum width=2.1cm, minimum height=0.75cm,
            font=\small\bfseries, text=primary] (C) at (5.2,0) {Prétraitement};
      \node[draw=primary, fill=lightblue, rounded corners=4pt,
            minimum width=2.1cm, minimum height=0.75cm,
            font=\small\bfseries, text=primary] (D) at (7.8,0) {Augmentation};
      \node[draw=primary, fill=lightblue, rounded corners=4pt,
            minimum width=2.1cm, minimum height=0.75cm,
            font=\small\bfseries, text=primary] (E) at (10.4,0) {Division};
      \draw[->, thick, color=secondary] (A.east) -- (B.west);
      \draw[->, thick, color=secondary] (B.east) -- (C.west);
      \draw[->, thick, color=secondary] (C.east) -- (D.west);
      \draw[->, thick, color=secondary] (D.east) -- (E.west);
    \end{tikzpicture}
    \vspace{1.5cm}
    {\small\color{darkgray}\textit{Pipeline complet — de la collecte à la préparation finale}}
  \end{center}
\end{titlepage}

\tableofcontents
\thispagestyle{fancy}
\newpage

%% ═══════════════════════════════════════════════════════════════════════════
\section{Introduction et Contexte}
%% ═══════════════════════════════════════════════════════════════════════════
Ce rapport documente l'intégralité du pipeline de constitution d'un dataset d'images 
d'oiseaux (5 espèces), depuis la collecte automatisée jusqu'à la préparation finale 
pour l'apprentissage profond.

\begin{infobox}[Objectif du projet]
Constituer un dataset propre, équilibré et annoté en respectant les bonnes pratiques 
de l'ingénierie des données et les principes éthiques de collecte.
\end{infobox}

\textbf{Technologies clés :}
\begin{itemize}[noitemsep]
  \item \textbf{Scraping :} Python, requests, BeautifulSoup, Selenium
  \item \textbf{Traitement image :} Pillow, OpenCV
  \item \textbf{Analyse :} pandas, matplotlib, scikit-learn
  \item \textbf{Versioning :} Git / GitHub (main + dev)
\end{itemize}

%% ═══════════════════════════════════════════════════════════════════════════
\section{Méthodologie Complète}
%% ═══════════════════════════════════════════════════════════════════════════

\subsection{Étape 1 — Collecte (Scraping)}
Le script \texttt{scraper\_oi.py} combine :
\begin{itemize}[noitemsep]
  \item \textbf{Requêtes statiques} (requests + BeautifulSoup) pour les pages simples.
  \item \textbf{Navigation dynamique} (Selenium) pour le chargement infini.
\end{itemize}
Un délai aléatoire (1–3 s) et une rotation d'User-Agent limitent les blocages.  
Les logs sont stockés dans \texttt{rapports/scraping.log}.

\subsection{Étape 2 — Nettoyage (\texttt{nettoyage.py})}
\begin{enumerate}[noitemsep]
  \item Suppression des doublons (hash MD5).
  \item Vérification d'intégrité (Pillow).
  \item Filtrage dimensionnel (rejet si < 64×64 px).
  \item Conservation des formats JPEG/PNG/WebP uniquement.
\end{enumerate}
Les opérations sont tracées dans \texttt{rapports/nettoyage.log}.

\subsection{Étape 3 — Prétraitement (\texttt{pretraitement.py})}
Standardisation de chaque image :
\begin{itemize}[noitemsep]
  \item Redimensionnement à 224×224 px (taille standard pour ResNet).
  \item Conversion RGB.
  \item Normalisation des pixels dans [0,1].
  \item Sauvegarde en tableau NumPy (.npy).
\end{itemize}

\subsection{Étape 4 — Augmentation (\texttt{augmentation.py})}
Pour enrichir la variabilité et corriger le déséquilibre initial :
\begin{itemize}[noitemsep]
  \item Rotation aléatoire (-15° à +15°)
  \item Symétrie horizontale (50\% proba)
  \item Zoom aléatoire (facteur 0.9–1.1)
  \item Luminosité / contraste ajustés (×0.8–1.2)
\end{itemize}
Chaque image originale génère 5 variantes.

\subsection{Étape 5 — Division (\texttt{division\_dataset.py})}
Répartition stratifiée (70\% train, 15\% val, 15\% test) avec graine aléatoire fixe (42).

%% ═══════════════════════════════════════════════════════════════════════════
\section{Difficultés Rencontrées et Solutions}
%% ═══════════════════════════════════════════════════════════════════════════
\begin{warnbox}[Principaux obstacles]
\begin{enumerate}[noitemsep]
  \item \textbf{Blocage HTTP (403/429)} → rotation d'User-Agent + délais adaptatifs.
  \item \textbf{Contenu JavaScript} → ajout de Selenium avec WebDriverWait.
  \item \textbf{Doublons inter-sources} → hachage MD5 global.
  \item \textbf{Déséquilibre de classes} → sur‑augmentation ciblée.
  \item \textbf{Images corrompues} → blocs try/except + log.
\end{enumerate}
\end{warnbox}
La pagination infinie a été gérée par simulation du défilement et détection de la fin de page via hauteur DOM.

%% ═══════════════════════════════════════════════════════════════════════════
\section{Statistiques Détaillées (avec Graphiques)}
%% ═══════════════════════════════════════════════════════════════════════════

\subsection{Bilan global du pipeline}
\begin{table}[H]
  \centering
  \caption{Tableau de synthèse complet}
  \renewcommand{\arraystretch}{1.3}
  \begin{tabular}{lc}
    \toprule
    \textbf{Étape} & \textbf{Nombre d'images} \\
    \midrule
    Collectées (scraping)              & 455 \\
    Supprimées (nettoyage)             &  62 (13,6\%) \\
    \quad dont filigranes              &  59 \\
    \quad dont doublons                &   3 \\
    Après nettoyage                    & 375 \\
    Après prétraitement                & 375 \\
    Après augmentation ($\times$6)     & 2 250 \\
    \midrule
    Train (70\%)  & 1 576 \\
    Validation (15\%) & 335 \\
    Test (15\%)   & 339 \\
    \bottomrule
  \end{tabular}
\end{table}

\subsection{Graphique 1 — Images collectées vs cible (100) par espèce}
\begin{center}
\begin{tikzpicture}
  \begin{axis}[
    ybar, bar width=18pt, width=13cm, height=7cm,
    symbolic x coords={Hibou,Flamant,Martin-p.,Cygne,Pic vert},
    xtick=data, ymin=0, ymax=115,
    ylabel={Nombre d'images}, xlabel={Espèce},
    nodes near coords, nodes near coords align={vertical},
    legend style={at={(0.98,0.98)}, anchor=north east},
    grid=major, grid style={dashed, gray!30},
    title={\textbf{Images collectées vs cible (100)}},
  ]
    \addplot[fill=primary!80] coordinates {
      (Hibou,87) (Flamant,100) (Martin-p.,100) (Cygne,68) (Pic vert,100)
    };
    \addplot[fill=none, draw=accent, very thick, dashed] coordinates {
      (Hibou,100) (Flamant,100) (Martin-p.,100) (Cygne,100) (Pic vert,100)
    };
    \legend{Collectées, Cible}
  \end{axis}
\end{tikzpicture}
\end{center}

\subsection{Graphique 2 — Suppressions par motif (camembert)}
\begin{center}
\begin{tikzpicture}
  \pie[
    radius=3.5, color={accent!80, warning!80, primary!80, success!80},
    text=pin, before number=\phantom{}, after number=\,\%
  ]{
    0/Corrompues,
    4.8/Taille insuffisante,
    95.2/Filigranes,
    0/Doublons
  }
\end{tikzpicture}
\captionof{figure}{Répartition des causes de suppression (dominées par les filigranes)}
\end{center}

\subsection{Graphique 3 — Avant / après augmentation}
\begin{center}
\begin{tikzpicture}
  \begin{axis}[
    ybar, bar width=14pt, width=13cm, height=7cm,
    symbolic x coords={Hibou,Flamant,Martin-p.,Cygne,Pic vert},
    xtick=data, ymin=0, ymax=620,
    ylabel={Nombre d'images}, legend style={at={(0.98,0.98)}, anchor=north east},
    grid=major, grid style={dashed, gray!30},
    nodes near coords, nodes near coords align={vertical},
    title={\textbf{Images finales vs images augmentées}},
  ]
    \addplot[fill=primary!80] coordinates {
      (Hibou,67)(Flamant,84)(Martin-p.,93)(Cygne,53)(Pic vert,78)
    };
    \addplot[fill=success!70] coordinates {
      (Hibou,402)(Flamant,504)(Martin-p.,558)(Cygne,318)(Pic vert,468)
    };
    \legend{Finales (après nettoyage), Augmentées ($\times$6)}
  \end{axis}
\end{tikzpicture}
\end{center}

\subsection{Graphique 4 — Répartition finale train/val/test (camembert global)}
\begin{center}
\begin{tikzpicture}
  \pie[radius=3, color={primary!80, secondary!80, accent!80}, text=pin, after number=\,\%]{
    70/Train (1 576),
    15/Validation (335),
    15/Test (339)
  }
\end{tikzpicture}
\captionof{figure}{Proportions respectées pour l'entraînement supervisé}
\end{center}

%% ═══════════════════════════════════════════════════════════════════════════
\section{Analyse Critique des Limites et Biais du Dataset}
%% ═══════════════════════════════════════════════════════════════════════════
\subsection{Biais géographique}
Les requêtes en français orientent Bing Images vers des sites européens.  
Le modèle risque d'être moins performant sur des oiseaux d'Afrique, d'Asie ou d'Amérique.

\subsection{Biais de saillance}
Les images les plus partagées (oiseaux colorés, poses typiques, plumage nuptial) sont surreprésentées.  
Les individus juvéniles, en mue ou partiellement cachés sont sous-représentés.

\subsection{Déséquilibre résiduel}
Même après augmentation, le ratio entre la classe la plus fréquente (Martin-pêcheur : 558) et la moins fréquente (Cygne : 318) est de 1,75.  
Des techniques de pondération ou de sur‑échantillonnage seront nécessaires en entraînement.

\subsection{Absence d'annotation humaine}
Les étiquettes sont déduites des noms de dossiers. Une vérification manuelle sur 5–10\% des images permettrait d'estimer le taux d'erreur.

\begin{warnbox}[Recommandations]
Ajouter un audit humain, documenter les dates de collecte, et réexaminer périodiquement la distribution.
\end{warnbox}

%% ═══════════════════════════════════════════════════════════════════════════
\section{Réflexion Éthique sur la Collecte d'Images}
%% ═══════════════════════════════════════════════════════════════════════════
\begin{ethicbox}[Principes appliqués]
Respect des \texttt{robots.txt}, délais inter‑requêtes, usage strictement académique.
\end{ethicbox}

\subsection{Droits d'auteur}
Les images appartiennent à leurs auteurs. Le dataset n'est pas redistribué publiquement et sert uniquement dans un cadre pédagogique (fair use).

\subsection{Consentement et vie privée}
Aucune photo de personne n'a été intentionnellement collectée. Si des visages apparaissaient accidentellement, ils seraient floutés.

\subsection{Biais algorithmique et justice}
Un modèle entraîné sur ce dataset pourrait être injuste envers des sous‑espèces ou des régions peu représentées. Cela doit être documenté dans toute publication.

\subsection{Transparence}
L'intégralité du code, des logs et de la méthodologie est accessible sur GitHub pour assurer la reproductibilité.

%% ═══════════════════════════════════════════════════════════════════════════
\section{Illustrations}
%% ═══════════════════════════════════════════════════════════════════════════


\vspace{1em}

\begin{figure}[H]
  \centering
  \begin{subfigure}{0.48\linewidth}
    \centering
    % Remplacer "screenshot1_scraping.png" par le nom réel de votre fichier
    \includegraphics[width=\linewidth, height=5cm, keepaspectratio]{1.png}
    \label{fig:scraping}
  \end{subfigure}
  \hfill
  \begin{subfigure}{0.48\linewidth}
    \centering
    \includegraphics[width=\linewidth, height=5cm, keepaspectratio]{2.png}
    \label{fig:nettoyage}
  \end{subfigure}
  \caption{Étapes 1 et 2 — Collecte et nettoyage}
  \label{fig:captures1}
\end{figure}

\begin{figure}[H]
  \centering
  \begin{subfigure}{0.48\linewidth}
    \centering
    \includegraphics[width=\linewidth, height=5cm, keepaspectratio]{3.png}
    \label{fig:augmentation}
  \end{subfigure}
  \hfill
  \begin{subfigure}{0.48\linewidth}
    \centering
    \includegraphics[width=\linewidth, height=5cm, keepaspectratio]{4.png}
    \label{fig:split}
  \end{subfigure}
  \caption{Étapes 4 et 5 — Augmentation et division du dataset}
  \label{fig:captures2}
\end{figure}

\begin{infobox}[Légende des captures]
\begin{enumerate}[noitemsep]
  \item \textbf{Fig. 1a} : Console du script de scraping — téléchargement des images par espèce.
  \item \textbf{Fig. 1b} : Résumé du nettoyage — images supprimées par motif (doublons, taille, filigranes).
  \item \textbf{Fig. 2a} : Visualisation des augmentations générées pour une image source.
  \item \textbf{Fig. 2b} : Tableau final de répartition des ensembles d'entraînement, validation et test.
\end{enumerate}
\end{infobox}


%% ═══════════════════════════════════════════════════════════════════════════
\section{Conclusion}
%% ═══════════════════════════════════════════════════════════════════════════
Ce projet a livré un pipeline complet, reproductible et documenté, produisant un dataset final de 2 250 images prêtes pour l'apprentissage.  
Les limites identifiées — biais géographique, de saillance, déséquilibre résiduel — sont clairement documentées, et des pistes d'amélioration sont proposées.

\begin{center}
\begin{tcolorbox}[colback=lightblue, colframe=primary, arc=6pt,
                  boxrule=1pt, width=0.75\linewidth]
  \centering
  \textbf{Rapport conforme aux livrables attendus}\\[0.3em]
  \small Scripts commentés, README, dataset organisé, analyse éthique et statistiques graphiques.
\end{tcolorbox}
\end{center}

\end{document}